% Actual degree-two cohomology, following the corrected S21.
% GPT-6 Astra Pro; /root/fixed_surface; 22 November 2026 UTC.
\section[HD-F1-S22: the degree-two excess map is injective]{HD-F1-S22: divisorial vanishing cycles determine the entire scalar degree-two cohomology of the actual norm fourfold}
\label{sec:f1-norm-fourfold-degree-two-support}

Retain the actual smooth canonical norm components $T_j$, the normal
Abel image $Y_j\subset P$, and the morphisms of S13--S21. Thus
$\pi:Z\to C$ is the specified connected etale cyclic triple cover,
$g(C)=5$, $g(Z)=13$, and
\[
 X=\operatorname{Sym}^8 Z,\qquad A=\operatorname{Pic}^8 C,
 \qquad F(D)=\mathcal O_C(\pi_*D),\qquad k=K_C.
 \tag{F1S22.1}
\]
The curve is in the same nonempty general-curve open used in
S13--S21; no additional hypothesis on its linear series is imposed
here. The conclusion is
\[
 H^2(P,\mathcal O_P)\xrightarrow{\ \sim\ }
 H^2(T_j,\mathcal O_{T_j})
   =H^2(Y_j,\mathcal O_{Y_j}),\qquad h^2(\mathcal O_{T_j})=28.
 \tag{F1S22.2}
\]
In particular the actual excess connecting map $\delta$ of
(F1S21.26) is injective. The proof below computes its kernel through
the topology and Hodge structures of the same morphism $F$.

\paragraph{1. The critical locus of the actual norm--Abel morphism.}
For an effective divisor $D=\sum_z m_z z$ on $Z$, define
\[
 E_{\max}(D)=\sum_{p\in C}\max_{z\in\pi^{-1}(p)}m_z\,p.
\]
The cotangent map of $F$ is the restriction of pulled-back
canonical forms:
\[
 (dF_D)^*:H^0(C,K_C)\longrightarrow H^0(Z,K_Z|_D),
 \qquad s\longmapsto\pi^*s|_D.
 \tag{F1S22.3}
\]
This follows from the cotangent map of the Abel morphism on a
symmetric product and the differential of the norm homomorphism.
It also proves the formula at repeated points: the etale local
coordinate at each point above $p$ identifies vanishing to order
$m_z$ with the corresponding jet condition on $s$ at $p$. Hence
\[
 \ker(dF_D)^*=H^0(C,K_C(-E_{\max}(D))).
 \tag{F1S22.4}
\]

Consider the closed incidence
\[
 I=\{(D,[s])\in X\times\mathbb P H^0(C,K_C):
                           D\le\pi^*\operatorname{div}(s)\}.
 \tag{F1S22.5}
\]
Its projection to $\mathbb P H^0(C,K_C)=\mathbb P^4$ is finite:
it is proper, and an effective divisor on a smooth curve has only
finitely many effective degree-eight subdivisors. The projection
to $X$ has fibre $\mathbb P\ker(dF_D)^*$. Consequently, if
$\mathrm{Crit}_r(F)$ denotes the locus where the cotangent kernel
has dimension at least $r$, then
\[
 \dim I=4,\qquad
 \dim\mathrm{Crit}_r(F)\le5-r\quad(r\ge1).
 \tag{F1S22.6}
\]
The equality $\dim I=4$ uses the existence of a degree-eight
subdivisor of every degree-twenty-four pulled-back canonical
divisor. In particular the entire critical locus has dimension
at most four, and the locus of corank at least two has dimension
at most three.

For $L\ne k$, Riemann--Roch gives $h^0(C,L)=4$.
The finite morphism $X\to\operatorname{Sym}^8 C$ is flat:
both symmetric products are smooth of dimension eight,
so the finite Cohen--Macaulay source over the regular
target satisfies miracle flatness. Its base change to
$|L|=\mathbb P^3$ is finite flat. The resulting fibre
is Cohen--Macaulay of pure dimension three. Thus every
fibre over $A\setminus\{k\}$ has pure dimension three.
The sole higher-dimensional fibre is the smooth fourfold union
\[
                         F^{-1}(k)=\mathcal T=\coprod_{j=0}^2T_j.
 \tag{F1S22.7}
\]
These statements concern the actual scheme fibres, as in S13
and S21.

\paragraph{2. A transverse slice at every generic divisorial critical value.}
Let $B$ be any irreducible divisor in the critical-value locus.
After deleting a proper closed subset of $B$, its inverse image
in the reduced critical locus is finite over $B$, every critical
point has rank four, and the restriction of $F$ to the critical
locus is etale onto $B$ at those points. Here finiteness follows
from (F1S22.6), while the last assertion follows from generic
smoothness in characteristic zero. Images of all exceptional
subsets have dimension at most three.

At such a point $y\in B$, the image of $dF_x$ at each critical
$x\in F^{-1}(y)$ is exactly $T_yB$: it contains the rank-four
image of the differential on the critical locus and has dimension
four. Choose local coordinates $(u_1,u_2,u_3,u_4,t)$ on $A$ with
$B=(t=0)$, and take the transverse disc $u_1=\cdots=u_4=0$.
Its inverse image $M$ is smooth of dimension four near the
whole central fibre. At a critical point this follows because
the four differentials $d(u_i\circ F)$ are independent; at
all other points it follows from smoothness of $F$.
After shrinking the disc, the proper map
\[
                    g:M\longrightarrow\Delta
 \tag{F1S22.8}
\]
is smooth away from the central fibre, and that fibre has only
finitely many singular points, each an isolated hypersurface
singularity of complex dimension three.

Write $M_0=F^{-1}(y)$ and $M_t=F^{-1}(y_t)$ for $t\ne0$.
The actual specialization triangle is
\[
 \mathbb Q_{M_0}\longrightarrow\psi_g\mathbb Q_M
 \longrightarrow\phi_g\mathbb Q_M\longrightarrow
 \mathbb Q_{M_0}[1].
 \tag{F1S22.9}
\]
The vanishing-cycle complex is supported at the finitely many singular
points and has cohomology only in degree three. One can obtain
this directly from Milnor's isolated hypersurface theorem.
Equivalently, $\mathbb Q_M[4]$ is perverse and
$\phi_g\mathbb Q_M[3]$ is perverse; a perverse sheaf supported
on finitely many points is concentrated in degree zero.
This second argument applies directly to the holomorphic
coordinates used in (F1S22.8).

Proper base change for nearby cycles and the long exact
sequence of (F1S22.9) give the actual specialization
isomorphisms
\[
 H^q(M_0,\mathbb Q)\xrightarrow{\ \sim\ }H^q(M_t,\mathbb Q)
               \quad\text{for }q=0,1,2,5,6.
 \tag{F1S22.10}
\]
They commute with monodromy; the source carries the identity
monodromy, so the monodromy on each displayed nearby group
is trivial. Degrees three and four are deliberately absent:
the nonzero degree-three vanishing cycles can affect both.

Choose a Whitney stratification of the proper map $F$ and
retain the open strata of each critical divisor just used.
There is a closed $Z_2\subset A$ of codimension at least two
such that on $A\setminus Z_2$ the only singular values are
these smooth divisorial strata. The transverse calculation,
together with local triviality along the strata, proves that
the local systems $\mathbb L_q=R^qF_*\mathbb Q$ from the
smooth locus extend over $A\setminus Z_2$ for $q=0,1,2,5,6$.
For the argument below it is enough to retain $q=0,1,2,6$.
The inclusion $A\setminus Z_2\hookrightarrow A$ induces an
isomorphism on fundamental groups: a loop, and a homotopy
between loops, can be moved off the real-codimension-at-least-four
strata by general position. Thus these local systems extend
uniquely to local systems $\overline{\mathbb L}_q$ on $A$.
This assertion about local systems does not assume any
base-change statement for coherent direct images.

\paragraph{3. The top cohomology already accounts for all divisorial components.}
On the smooth locus, relative Serre duality and (F1S21.18)
give
\[
 h^0(F^{-1}(y),\mathcal O)=3,\qquad
 h^2(F^{-1}(y),\mathcal O)=84.
 \tag{F1S22.11}
\]
Indeed $R^3F_*\omega_X$ has rank three and
$R^1F_*\omega_X$ has rank eighty-four; the relative
canonical line equals $\omega_X$ after trivializing
$\omega_A$. A smooth projective fibre therefore has three
connected components. Poincare duality in dimension three
identifies
\[
        \mathbb L_6\simeq\mathbb L_0^\vee(-3),
                  \qquad\operatorname{rank}\mathbb L_6=3.
 \tag{F1S22.12}
\]
This also identifies their extensions on $A$.

At a general point of any critical divisor, (F1S22.10)
implies $\dim H^6(F^{-1}(y),\mathbb Q)=3$.
The same holds at a general point of a divisor outside the
critical-value locus by smoothness. Since a proper pure
three-dimensional complex variety has one rational top
cohomology generator for each irreducible component, this
proves that a generic divisorial fibre has exactly three
irreducible components. More precisely, all of its top
cohomology is already supplied by the extending local system
$\overline{\mathbb L}_6$; the next paragraph uses this
equality inside the actual direct-image decomposition.

\paragraph{4. Every support and shift that can contribute in degree two.}
Put
\[
                     \mathcal H=RF_*\mathbb Q_X[8].
 \tag{F1S22.13}
\]
The projective decomposition theorem and relative hard
Lefschetz give
\[
 \mathcal H\simeq\bigoplus_i{}^pH^i(\mathcal H)[-i],
 \qquad
 \ell^i:{}^pH^{-i}(\mathcal H)
                     \xrightarrow{\ \sim\ }{}^pH^i(\mathcal H),
 \tag{F1S22.14}
\]
with the Tate twists inserted in the Hodge-module version.
Each perverse cohomology is a direct sum of intersection
complexes $IC_B(\mathbb V)$. Our convention is that on a
smooth dense open subset of a support of dimension $d$,
$IC_B(\mathbb V)=\mathbb V[d]$. Moreover
\[
             \mathcal H^a(IC_B(\mathbb V))=0\quad(a<-d).
 \tag{F1S22.15}
\]
Consequently a term $IC_B(\mathbb V)[-i]$ contributes at
a generic point of $B$ to fibre cohomology in degree
$8+i-d$, and at any point of $B$ only in degrees at
least $8+i-d$.

Suppose first that $B\subsetneq A$ has codimension $b$
and $B\ne\{k\}$. Its generic fibre has dimension three,
so nonzero generic stalk cohomology forces
$8+i-(5-b)\le6$, that is $i\le3-b$.
Relative hard Lefschetz preserves the strict supports
and pairs the support in index $i$ with the same support
in index $-i$. Applying the same bound to that index gives
\[
                   b-3\le i\le3-b.
 \tag{F1S22.16}
\]
Together with (F1S22.15), this says that its contribution
to any fibre occurs only in degrees at least
\[
                         8+i-(5-b)\ge2b.
 \tag{F1S22.17}
\]
Thus supports of codimension at least two do not
contribute to degree two. If the support is the point
$\{k\}$ itself, its four-dimensional fibre gives $i\le0$,
and relative hard Lefschetz gives $i\ge0$. Its only
possible index is zero, contributing in degree eight.

For a divisorial support $B$, (F1S22.16) permits
$-2\le i\le2$. The only term that can contribute in
degree two is $IC_B(\mathbb V)[2]$, namely $i=-2$.
Relative hard Lefschetz would give a nonzero companion
$IC_B(\mathbb V')[-2]$ in index $i=2$. At a general point
of $B$ that companion contributes to $H^6(F^{-1}(y),\mathbb Q)$.
But the full-support term from $\mathbb L_6$ is
\[
          IC_A(\mathbb L_6)[-3]
                 =\overline{\mathbb L}_6[2],
 \tag{F1S22.18}
\]
and already contributes dimension three there. By paragraph
3 this is the entire dimension of $H^6$.
The direct-sum decomposition therefore leaves no room
for the proposed nonzero companion. This excludes every
divisorial term with $i=2$ and, by the actual Lefschetz
isomorphism, every divisorial term with $i=-2$.

It remains to check the full-support terms. On the smooth
locus, $\mathbb L_q$ occurs in perverse index $i=q-3$.
For $q=0,1,2$ its intermediate extension equals
$\overline{\mathbb L}_q[5]$, so the shifted term has
stalk only in degree $q-8$ of $\mathcal H$.
For $q\ge3$, (F1S22.15) bounds the stalk degrees below
by $-5+(q-3)=q-8\ge-5$.
It follows that the entire degree-minus-six stalk of
$\mathcal H$, equivalently $H^2(\mathcal T,\mathbb Q)$,
comes precisely from
\[
       IC_A(\mathbb L_2)[1]
                 =\overline{\mathbb L}_2[6].
 \tag{F1S22.19}
\]

\paragraph{5. Hodge compatibility and the actual restriction isomorphism.}
Use the decomposition in (F1S22.14) in the derived
category of mixed Hodge modules, starting from the
constant pure Hodge module $\mathbb Q_X^H[8]$.
Its perverse direct images are pure of weight $8+i$,
and have the same strict-support decomposition as their
underlying rational perverse sheaves. The full-support
summand corresponding to $\mathbb L_2$ has weight seven.
Its underlying rational complex is
$\overline{\mathbb L}_2[5]$, a shifted local system on
the smooth variety $A$. It is therefore a smooth pure
Hodge module, equivalently a polarizable variation of
Hodge structure of weight $7-5=2$ on all of $A$.

Apply the actual derived pullback to the point $k$ in
the category of mixed Hodge modules and then take
cohomology in degree minus six. Proper base change
identifies the result on the left with the usual Hodge
structure on $H^2(\mathcal T,\mathbb Q)$.
The support and shift calculation proves that every
summand except (F1S22.19) has zero underlying rational
cohomology in that degree, hence zero Hodge object.
We obtain an isomorphism of Hodge structures
\[
 H^2(\mathcal T,\mathbb Q)
          \simeq (\overline{\mathbb L}_2)_k.
 \tag{F1S22.20}
\]
No inference from equality of Betti numbers alone is
being used. The Hodge numbers of the variation on the
right are constant and agree with those of a smooth
generic fibre. Equations (F1S22.11) and (F1S22.20) give
\[
                   h^2(\mathcal T,\mathcal O_{\mathcal T})=84.
 \tag{F1S22.21}
\]
The actual deck transformation permutes the three $T_j$
transitively, so each has $h^2(\mathcal O)=28$.
S18 already proves that the actual restriction map
$H^2(P,\mathcal O_P)\to H^2(T_j,\mathcal O_{T_j})$
has rank twenty-eight. This identifies the isomorphism
in (F1S22.2) with the prescribed restriction map, rather
than only with an abstract equality of dimensions.
S16--S17 identify $Rb_*\mathcal O_{T_j}$ with
$\mathcal O_{Y_j}$, giving the stated equality for $Y_j$.

Finally the exact sequence (F1S21.26) is
\[
 0\longrightarrow H^2(X,\mathcal K)
 \longrightarrow H^2(\mathcal T,\mathcal O_{\mathcal T})
 \longrightarrow H^4(\mathcal T,\mathcal O_{\mathcal T}(-h))
 \xrightarrow{\delta}H^3(X,\mathcal K),
 \tag{F1S22.22}
\]
where $\mathcal K=LF^*k_k$ and the first term has
dimension eighty-four. Its injection is the one
constructed in S21; (F1S22.21) makes it an isomorphism.
Exactness therefore proves $\ker\delta=0$.

\paragraph{6. The precise consequence for the full trace.}
The actual scalar morphism of S18,
\[
 H^2(Y_j,\mathcal O_{Y_j})\longrightarrow
 \operatorname{Ext}^2_P(\mathcal O_{Y_j},\mathcal O_{Y_j}),
 \tag{F1S22.23}
\]
is a split injection supplied by multiplication and the
unit on $\mathcal O_{Y_j}\otimes^L_{\mathcal O_P}\mathcal O_{Y_j}$.
Now (F1S22.2) identifies its entire source with the
twenty-eight ambient scalar classes checked in S18.
Thus the full Buchweitz--Flenner trace is injective on
the whole scalar degree-two subspace for this actual
object. The calculation of its kernel on the rest of
$\operatorname{Ext}^2$ remains a separate, explicit
self-Ext computation; no dimension for that rest is
deduced here.

\paragraph{Primary inputs and proof scope.}
The perverse decomposition, strict-support decomposition,
relative hard Lefschetz, and the ordinary-degree bound
(F1S22.15) are the results of de Cataldo--Migliorini,
\emph{The Hodge theory of algebraic maps}, Annales ENS
38 (2005), Theorem~2.1.1 and Section~3.8, with their
projective hypotheses satisfied by the displayed $F$.
For Hodge compatibility we use Saito's projective direct
image theorem for polarizable Hodge modules,
\emph{Modules de Hodge polarisables}, PRIMS~24 (1988),
Theorem~5.3.1 and its complete Section~5.3 proof,
pp.~977--988, and the derived functors in
\emph{Mixed Hodge Modules}, PRIMS~26 (1990),
Sections~4.3--4.4, in particular the base-change
isomorphisms (4.4.3) and the pure decomposition (4.5.4).
The exact smooth-object
criterion used in paragraph~5 is Theorem~3.27 and its
algebraic remark on p.~313 of that same original paper.
The complete Section~5.3 proof in the 1988 paper and
the relevant proof of Theorem~3.27, including the adjoining
normal-crossing propositions, in the 1990 paper were read
in the publisher texts. The direct-image mechanism is
also described in Saito's \emph{A young person's guide
to mixed Hodge modules}, Theorems~1.3--1.4 and their
proof discussion, and Schnell's complete Section~17
account of the projective direct-image argument.
The isolated-singularity calculation is Milnor,
\emph{Singular Points of Complex Hypersurfaces},
Chapter~6, or directly the perversity statement and
actual Milnor-fibre identification in Dimca--Saito,
\emph{Some consequences of perversity of vanishing
cycles}, Annales de l'Institut Fourier~54 (2004),
1769--1792, DOI~10.5802/aif.2065, Section~1.1
(with the nearby-cycle maps in Sections~1.2--1.3).
The application uses the actual
specialization triangle, not a presumed splitting of
the excess Koszul complex in S21.
